\documentclass{article}

\usepackage{amsmath,amssymb}
\usepackage{xurl}
\usepackage[colorlinks=true,linkcolor=blue,citecolor=blue,urlcolor=blue]{hyperref}
\setlength{\emergencystretch}{2em}

% Shared commands for notes in docs/paper-gaps/.

\DeclareUnicodeCharacter{2080}{\ifmmode{}_{0}\else\textsubscript{0}\fi}
\DeclareUnicodeCharacter{2081}{\ifmmode{}_{1}\else\textsubscript{1}\fi}
\DeclareUnicodeCharacter{2082}{\ifmmode{}_{2}\else\textsubscript{2}\fi}
\DeclareUnicodeCharacter{2099}{\ifmmode{}_{n}\else\textsubscript{n}\fi}

\newcommand{\Lean}{\textsc{Lean}}
\ExplSyntaxOn
\NewDocumentCommand{\leanid}{m}{
  \ifmmode
    \textsf{\detokenize{#1}}
  \else
    \group_begin:
    \urlstyle{sf}
    \tl_set:Nn \l_tmpa_tl {#1}
    \tl_replace_all:Nnn \l_tmpa_tl {\_} {_}
    \exp_args:NV \path \l_tmpa_tl
    \group_end:
  \fi
}
\ExplSyntaxOff
\providecommand{\tr}{\operatorname{tr}}
\newcommand{\tnleanIssueBase}{https://github.com/LionSR/TNLean/issues/}
\newcommand{\tnleanPrBase}{https://github.com/LionSR/TNLean/pull/}
\newcommand{\tnleanDocBase}{https://sirui-lu.com/TNLean/docs/}
\newcommand{\ghissue}[1]{\href{\tnleanIssueBase#1}{GitHub issue \##1}}
\newcommand{\ghpr}[1]{\href{\tnleanPrBase#1}{PR \##1}}
\newcommand{\doclink}[1]{\href{\tnleanDocBase#1.html}{\path{#1}}}

% Dirac notation and the identity operator, as in blueprint/src/macros/common.tex.
% \providecommand keeps note-local definitions authoritative.
\usepackage{dsfont}
\providecommand{\ket}[1]{|#1\rangle}
\providecommand{\bra}[1]{\langle #1|}
\providecommand{\braket}[2]{\langle #1 | #2 \rangle}
\providecommand{\ketbra}[2]{|#1\rangle\!\langle #2|}
\providecommand{\Id}{\mathds{1}}
\providecommand{\one}{\mathds{1}}
\providecommand{\C}{\mathbb{C}}
\providecommand{\MN}[1]{M_{#1}(\C)}
\providecommand{\MD}{M_D(\C)}

% External referencing for paper sources.  The build helper writes sanitized
% auxiliary files under build/paper-aux/ and excludes duplicated source labels.
\usepackage{xr}

% Import a generated paper auxiliary file with a caller-supplied label prefix.
% The candidate paths support builds from docs/paper-gaps/, the repository root,
% and build/paper-aux/ itself.
\newcommand{\importpaperaux}[2]{%
  \IfFileExists{../../build/paper-aux/#2.aux}{%
    \externaldocument[#1]{../../build/paper-aux/#2}%
  }{%
    \IfFileExists{build/paper-aux/#2.aux}{%
      \externaldocument[#1]{build/paper-aux/#2}%
    }{%
      \IfFileExists{#2.aux}{%
        \externaldocument[#1]{#2}%
      }{}%
    }%
  }%
}

% General external reference macros
\newcommand{\paperref}[2]{\ref{#1-#2}}
\newcommand{\papereqref}[2]{(\ref{#1-#2})}

% CPSV16 source (1606.00608 / MPDO-22-12-17-2.tex) external document and shortcuts
\importpaperaux{cpsv16-}{1606.00608/MPDO-22-12-17-2-tnlean}
\newcommand{\cpsvref}[1]{\paperref{cpsv16}{#1}}
\newcommand{\cpsveqref}[1]{\papereqref{cpsv16}{#1}}

% Verdict marker: \gapnote{<kind>}{<status>}. Kind is the policy
% classification (clarification, local-correction, scope-restriction,
% unfaithful, false-source, open-gap); status is open, wip (elimination
% underway), resolved, or historical. Machine-read by the site index;
% prints a verdict line.
\newcommand{\gapnote}[2]{\par\noindent\textbf{Verdict:} #1 (#2).\par}


\title{Canonical-Form Weight Normalization and the Proportional MPV Comparison}
\author{}
\date{2026-05-10}

\begin{document}
\maketitle
\gapnote{unfaithful}{open}

\section{The source argument}

Cirac--P\'erez-Garc\'ia--Schuch--Verstraete build their multi-block proof of
the Fundamental Theorem on a canonical form
\cite[Section~\cpsvref{Sec:MPV}]{Cirac2016MPDO_arXiv}. Let $A=(A^i)_{i=1}^d$ be
a tensor with physical dimension $d$. In the source canonical form, a
bond-space basis change gives
\begin{align}
    A^i
        &= \bigoplus_{k=1}^r \mu_k A_k^i,
    \label{eq:cpsv16_cf_normalization_and_proportional_comparison_source_canonical_form}
\end{align}
where each block tensor $A_k$ is normal and each scalar $\mu_k$ is nonzero.
This is Equation~(\cpsvref{eq:II_CF1}) of Ref.~\cite{Cirac2016MPDO_arXiv}.
Immediately after that equation, the source sets the convention
\begin{quote}
    ``we can always choose $|\mu_k|\leq 1$ and at least one of them equals one,
    something which we will assume from now on.''
    \cite[paragraph after Eq.~(\cpsvref{eq:II_CF1})]{Cirac2016MPDO_arXiv}
\end{quote}
This normalization is part of the working definition throughout the paper
rather than a downstream consequence. Under the strict modulus ordering of the
basis of normal tensors, the dominant block satisfies $|\mu_1|=1$, while every
sub-dominant block satisfies $|\mu_k|<1$.

The same paper proves the proportional Fundamental Theorem and its equal-MPV
corollary \cite[Section~\cpsvref{Sec:MPV}]{Cirac2016MPDO_arXiv}. Consider
tensors $A$ and $B$ written in basis-of-normal-tensors form. Let $(A_j)_{j=1}^g$
and $(B_j)_{j=1}^g$ be their representative families, and let $\mu_{j,q_a}$
and $\nu_{j,q_b}$ be the corresponding coefficients, where
$q_a=1,\ldots,r_{a,j}$ and $q_b=1,\ldots,r_{b,j}$. The proportional Fundamental
Theorem identifies the representatives by
\begin{align}
    B_j
        &= e^{i\phi_j}Y_j A_j Y_j^{-1},
    \label{eq:cpsv16_cf_normalization_and_proportional_comparison_representative_gauge}
\end{align}
for phases $\phi_j$ and invertible matrices $Y_j$. The trace-pair argument in
the proof of the equal-MPV corollary then gives the scalar identity
\begin{align}
    \sum_{q=1}^{r_{a,j}}\mu_{j,q}^N
        &= \sum_{q=1}^{r_{b,j}}
            (\nu_{j,q}e^{i\phi_j})^N.
    \label{eq:cpsv16_cf_normalization_and_proportional_comparison_source_power_sum}
\end{align}
This is the unlabeled identity used in the source proof
\cite[Section~\cpsvref{Sec:MPV}]{Cirac2016MPDO_arXiv}. The elementary power-sum
lemma cited there gives
\begin{align}
    r_{a,j}
        &= r_{b,j}
        =: r_j,
    \label{eq:cpsv16_cf_normalization_and_proportional_comparison_equal_multiplicities}\\
    \{\mu_{j,q}:1\leq q\leq r_j\}
        &= \{\nu_{j,q}e^{i\phi_j}:1\leq q\leq r_j\},
    \label{eq:cpsv16_cf_normalization_and_proportional_comparison_weight_multiset}
\end{align}
where the second equality holds for multisets. Matching bond dimensions and
the global gauge formula follow.

The source extracts coefficient data directly from the power sums in
Eq.~\ref{eq:cpsv16_cf_normalization_and_proportional_comparison_source_power_sum}.
It feeds a finite set of lengths $N$ into a separation lemma for sums of $N$th
powers, without invoking the convergence of any quantity indexed by $N$.

\section{The formal canonical form}

The formal counterpart of the canonical form in the basis-of-normal-tensors
route is the normal canonical form structure for BNT representations.\footnote{%
    Formal structure: \leanid{IsNormalCanonicalFormBNT} in
    \path{TNLean/MPS/BNT/Construction.lean}, lines 181--189.}
Its data consist of a weight family $\mu:\mathrm{Fin}\,r\to\C$ with strict
modulus ordering, nonzero weights $\mu_k\neq 0$ for each $k$, and the
blockwise normality inherited from the underlying canonical form. In the version
examined here, the modulus of the dominant weight was unconstrained.

The source definition additionally requires
\begin{align}
    |\mu_1|
        &= 1,
    \label{eq:cpsv16_cf_normalization_and_proportional_comparison_missing_dominant_normalization}
\end{align}
where the source indexes from $1$ and the formalization indexes from $0$. The
blueprint chapter on the basis of normal tensors records strict modulus
ordering across its separated canonical-form definitions. The canonical-form
chapter separately records a phase--modulus decomposition
\begin{align}
    \mu_k
        &= \eta_k|\mu_k|,
    \label{eq:cpsv16_cf_normalization_and_proportional_comparison_phase_modulus_decomposition}
\end{align}
with $\eta_k$ on the unit circle. Neither location originally imposed
Eq.~\ref{eq:cpsv16_cf_normalization_and_proportional_comparison_missing_dominant_normalization}.\footnote{%
    The basis-of-normal-tensors definitions are
    \texttt{def:cf\_bnt} and \texttt{def:normal\_cf\_bnt} in
    \path{blueprint/src/chapter/ch10_bnt.tex}, lines 70--100 and 119--141,
    respectively. The phase--modulus decomposition is recorded in
    \path{blueprint/src/chapter/ch09_canonical.tex}, lines 125--139.}

This omission affects downstream arguments. The source normalization places the
dominant block on the unit circle, as required by spectral and irreducibility
arguments, while placing every sub-dominant block strictly inside the open unit
disk. The latter property drives the comparison with the discarded formal limit
hypothesis.

\section{The discarded formal proportional comparison}

The discarded formal route to the proportional Fundamental Theorem replaced
the source power-sum identity in
Eq.~\ref{eq:cpsv16_cf_normalization_and_proportional_comparison_source_power_sum}
with coefficient sequences indexed by chain length. For each block index $j$,
it introduced a sequence $N\mapsto a_{N,j}\in\C$ and a designated limit
$a_j^\infty\in\C$ satisfying
\begin{align}
    a_{N,j}
        &\xrightarrow[N\to\infty]{} a_j^\infty,
    \label{eq:cpsv16_cf_normalization_and_proportional_comparison_formal_coefficient_limit}\\
    a_j^\infty
        &\neq 0.
    \label{eq:cpsv16_cf_normalization_and_proportional_comparison_formal_nonzero_limit}
\end{align}
It imposed analogous data $b_{N,k}\to b_k^\infty\neq 0$ on the $B$-side, along
with proportionality scalars $c_N$ that approach a nonzero limit $c$:
\begin{align}
    c_N
        &\xrightarrow[N\to\infty]{} c,
    \label{eq:cpsv16_cf_normalization_and_proportional_comparison_scalar_limit}\\
    c
        &\neq 0.
    \label{eq:cpsv16_cf_normalization_and_proportional_comparison_nonzero_scalar_limit}
\end{align}
This route identified the scalars by comparing limits rather than by applying
the source power-sum lemma.

Consider the natural specialization where the $j$th source sector contains a
single representative, so that $r_j=1$ in multiplicity notation. Letting
$\mu_j^A$ denote the corresponding source weight, the formal coefficient array
was set to
\begin{align}
    a_{N,j}
        &= (\mu_j^A)^N.
    \label{eq:cpsv16_cf_normalization_and_proportional_comparison_natural_coefficient_array}
\end{align}
The associated blueprint passage appears in remark
\texttt{rem:cf\_coeff\_arrays} following \texttt{lem:ft\_proportional} in
\path{blueprint/src/chapter/ch11_fundamental_theorem_proof.tex},
lines 172--202. Under
Eq.~\ref{eq:cpsv16_cf_normalization_and_proportional_comparison_natural_coefficient_array},
the requirements in
Eqs.~\ref{eq:cpsv16_cf_normalization_and_proportional_comparison_formal_coefficient_limit}
and
\ref{eq:cpsv16_cf_normalization_and_proportional_comparison_formal_nonzero_limit}
become conditions on the powers $(\mu_j^A)^N$.

\section{The design history of the limit hypothesis}

The convergence-to-nonzero-limit hypothesis appears neither in
Ref.~\cite{Cirac2016MPDO_arXiv} nor in the initial proof plan for the
formalization. A memo committed on 2026-02-22 outlines the planned proof of the
proportional Fundamental Theorem and its equal-MPV corollary from that
source.\footnote{%
    \path{docs/audits/2026-02-22_thm4_4_proof_plan_ref1606.txt}, committed in
    \texttt{64c018ab} earlier the same day as the introducing implementation
    \texttt{dd54b79e}.}
That memo proposes a span-equality bridge alongside a finite-length scalar
power-sum identity to match the $\mu$-coefficients.

In its own notation, the memo records the scalar identity
\begin{align}
    \sum_{k=1}^{x_a}\lambda_{a,k}^N
        &= \sum_{k=1}^{x_b}\lambda_{b,k}^N.
    \label{eq:cpsv16_cf_normalization_and_proportional_comparison_memo_scalar_power_sum}
\end{align}
It proposes formalizing
Eq.~\ref{eq:cpsv16_cf_normalization_and_proportional_comparison_memo_scalar_power_sum}
as a fixed-length theorem evaluated at finitely many values of $N$. The planned
basis-of-normal-tensors application is
\begin{align}
    \sum_{q=1}^{r_{a,j}}\mu_{j,q}^N
        &= \sum_{q=1}^{r_{b,j}}
            (\nu_{j,q}e^{i\phi_j})^N.
    \label{eq:cpsv16_cf_normalization_and_proportional_comparison_memo_bnt_power_sum}
\end{align}
This scalar theorem would then establish
Eqs.~\ref{eq:cpsv16_cf_normalization_and_proportional_comparison_equal_multiplicities}
and
\ref{eq:cpsv16_cf_normalization_and_proportional_comparison_weight_multiset}.
The memo contains neither a convergence statement nor a nonzero-limit
hypothesis.

The memo also identifies an apparent circularity in the span-equality route.
Let $V^{(N)}(A_j)$ denote the length-$N$ periodic matrix product vector
generated by $A_j$, with $V^{(N)}(B_k)$ defined analogously. Deducing span
equality from proportionality seems to require linear independence at a witness
length, while the surrounding basis-of-normal-tensors permutation step uses
span equality to obtain that linear independence. The memo suggests replacing
global span equality with eventual span equality:
\begin{align}
    \operatorname{span}\{V^{(N)}(A_j)\}_j
        &= \operatorname{span}\{V^{(N)}(B_k)\}_k
        \qquad\text{for all sufficiently large }N.
    \label{eq:cpsv16_cf_normalization_and_proportional_comparison_eventual_span_equality}
\end{align}
Because the surrounding argument needs only this eventual form, the proposed
weakening resolves the circularity at the type level.

The implementation committed later that day took neither route. It omitted the
planned fixed-length scalar identity for the $\mu$-coefficient comparison and
set aside
Eq.~\ref{eq:cpsv16_cf_normalization_and_proportional_comparison_eventual_span_equality}.
Instead, it added global convergence and nonzero-limit hypotheses to the
basis-of-normal-tensors bridge theorems, later bundling them into a single data
structure. These fields reflect neither the source power-sum identity nor the
memo's eventual span equality.

The project record gives no contemporaneous mathematical justification for
this third approach. A commit on 2026-05-09 describes the resulting structure
as ``MPV-level Newton--Girard convergence data, not a \S II 283--302
Wedderburn-style similarity.'' That comment describes the resulting structure,
but it does not explain why the implementation diverged from the
memo.\footnote{%
    The introducing commit is \texttt{dd54b79e}, ``Add ScalarPowerSumIdentity,
    BNTConstruction, BNTPermutationThm44.'' Despite the name, the
    \leanid{ScalarPowerSumIdentity} module added by that commit does not
    implement the fixed-length scalar identity planned from the elementary
    power-sum lemma in Ref.~\cite{Cirac2016MPDO_arXiv}. The structure refactor
    that bundled the limit fields is \texttt{9e096a14} (2026-03-22), and the
    later Newton--Girard description appears in the message of
    \texttt{8ab1a5db} (2026-05-09).}

The design history thus confirms that the limit hypothesis comes from neither
the cited source nor the recorded proof plan.

\section{The mathematical incompatibility}

This section focuses specifically on the strict-weight specialization with one
representative per block:
\begin{align}
    a_{N,j}
        &= (\mu_j^A)^N,
    \label{eq:cpsv16_cf_normalization_and_proportional_comparison_specialized_a_array}\\
    b_{N,k}
        &= (\mu_k^B)^N.
    \label{eq:cpsv16_cf_normalization_and_proportional_comparison_specialized_b_array}
\end{align}
It does not address the broader sector coefficient arrays
\begin{align}
    c_{S,j}^{(N)}
        &= \sum_q \mu_{j,q}^N,
    \label{eq:cpsv16_cf_normalization_and_proportional_comparison_sector_coefficient_array}
\end{align}
which retain repeated copies within a basis-of-normal-tensors sector.

Under the source normalization,
\begin{align}
    |\mu_1^A|
        &= 1,
    \label{eq:cpsv16_cf_normalization_and_proportional_comparison_source_dominant_modulus}\\
    |\mu_j^A|
        &< 1
        \qquad (j\geq 2).
    \label{eq:cpsv16_cf_normalization_and_proportional_comparison_source_subdominant_modulus}
\end{align}
For every sub-dominant block,
\begin{align}
    |(\mu_j^A)^N|
        &= |\mu_j^A|^N,
    \label{eq:cpsv16_cf_normalization_and_proportional_comparison_power_modulus}\\
    |\mu_j^A|^N
        &\xrightarrow[N\to\infty]{} 0.
    \label{eq:cpsv16_cf_normalization_and_proportional_comparison_subdominant_decay}
\end{align}
The only possible limit of $a_{N,j}$ is therefore
\begin{align}
    a_j^\infty
        &= 0
        \qquad (j\geq 2),
    \label{eq:cpsv16_cf_normalization_and_proportional_comparison_subdominant_zero_limit}
\end{align}
which contradicts
Eq.~\ref{eq:cpsv16_cf_normalization_and_proportional_comparison_formal_nonzero_limit}.

For the dominant block, the modulus of $(\mu_1^A)^N$ remains $1$. The sequence
converges only when $\mu_1^A=1$; a nontrivial root of unity yields a periodic,
nonconvergent sequence, while a generic phase also fails to converge.

The mismatch is structural. Under the source normalization, a multi-block
family with $r>1$ cannot instantiate the discarded data through
Eq.~\ref{eq:cpsv16_cf_normalization_and_proportional_comparison_specialized_a_array},
because each sub-dominant block violates the nonzero-limit requirement.
Dropping the source normalization might make the formal limit hypothesis
satisfiable, but the resulting tensor family no longer satisfies the
canonical-form conditions of Ref.~\cite{Cirac2016MPDO_arXiv}.

The source proof avoids this obstruction because it never takes the large-$N$
limit of $(\mu_j^A)^N$. Instead, it recovers the weight multiset from a finite
collection of power-sum identities in $N$. In the strict one-representative
setting, once basis-of-normal-tensors linear independence provides a witness
length $N_0$, evaluating at $N_0$ fixes the coefficients up to the standard
phase or root-of-unity factor. In the general setting with multiplicities, the
source separation lemma and the 2026-02-22 memo use the finite range
$N=1,\ldots,\max(r_{a,j},r_{b,j})$.

\section{Cross-check with the blueprint}

The blueprint entries that once mirrored the convergence and nonzero-limit
interface have been demoted or removed from the source-aligned theorem path.
They should not be treated as formalizations of the proportional Fundamental
Theorem or its equal-MPV corollary from Ref.~\cite{Cirac2016MPDO_arXiv}. The
active blueprint now tracks the source comparison as an open block-matching
step that relies on the basis-of-normal-tensors power-sum argument rather than
external coefficient limits.

The blueprint and the original formal definition agreed on one point: both
recorded strict modulus ordering while omitting
Eq.~\ref{eq:cpsv16_cf_normalization_and_proportional_comparison_missing_dominant_normalization}.

\section{What a corrected formalization would say}

A source-faithful formalization requires two adjustments.

First, the canonical-form data must include
Eq.~\ref{eq:cpsv16_cf_normalization_and_proportional_comparison_missing_dominant_normalization}.
Strict modulus ordering then ensures
\begin{align}
    |\mu_k|
        &\leq 1
        \qquad (1\leq k\leq r),
    \label{eq:cpsv16_cf_normalization_and_proportional_comparison_all_weight_bound}
\end{align}
with strict inequality on every sub-dominant block.\footnote{%
    The formal field carrying the strict ordering is
    \leanid{mu_strict_anti} in
    \leanid{IsNormalCanonicalFormBNT}; see
    \path{TNLean/MPS/BNT/Construction.lean}, lines 184--185.}
This dominant normalization belongs in the formal BNT canonical form structure
and the corresponding blueprint definition.\footnote{Formal structure:
\leanid{IsNormalCanonicalFormBNT}.} It formalizes the source convention
directly rather than imposing an additional assumption.

Second, the proportional comparison must use finite power-sum relations. In the
general setting with coefficient multiplicities, the blockwise hypothesis takes
the form
\begin{align}
    \sum_{q=1}^{r_{a,j}}\mu_{j,q}^N
        &= \sum_{q=1}^{r_{b,j}}
            (\nu_{j,q}e^{i\phi_j})^N,
    \label{eq:cpsv16_cf_normalization_and_proportional_comparison_corrected_finite_power_sum}\\
    1
        &\leq N
        \leq \max(r_{a,j},r_{b,j}).
    \label{eq:cpsv16_cf_normalization_and_proportional_comparison_finite_power_sum_range}
\end{align}
Equivalently, one may evaluate at the finite set of witness lengths supplied
once span and linear-independence hypotheses are available. The source
separation lemma then yields
Eqs.~\ref{eq:cpsv16_cf_normalization_and_proportional_comparison_equal_multiplicities}
and
\ref{eq:cpsv16_cf_normalization_and_proportional_comparison_weight_multiset}.
This matches the route proposed in the 2026-02-22 memo for arbitrary
multiplicities, where a single witness length does not suffice.

In the strict-weight specialization with a single representative, the
comparison reduces to a single witness length $N_0$. Linear independence of the
basis of normal tensors at length $N_0$ comes from the direct-sum span theorem.
For a proportionality scalar $c$, the blockwise relation becomes
\begin{align}
    (\mu_j^A)^{N_0}
        &= c^{N_0}(\mu_j^B)^{N_0},
        \qquad j=1,\ldots,g.
    \label{eq:cpsv16_cf_normalization_and_proportional_comparison_single_witness_weights}
\end{align}
This identity determines $\mu_j^A$ from $\mu_j^B$ and $c$ up to an $N_0$th
root of unity, an ambiguity absorbed by the phase $\phi_j$ in
Eq.~\ref{eq:cpsv16_cf_normalization_and_proportional_comparison_representative_gauge}.
The argument requires no convergence hypothesis or nonzero limit for
$(\mu_j^A)^N$.

The corrected route is finite and follows the source. The multi-multiplicity
setting relies on the finite power-sum identities from the proof of the
equal-MPV corollary \cite[Section~\cpsvref{Sec:MPV}]{Cirac2016MPDO_arXiv}, with
the single-witness argument serving as its strict-weight specialization. With
linear independence provided by the direct-sum span theorem, this approach also
bypasses the circularity noted in the 2026-02-22 memo for the span-equality
bridge.

\section{Conclusion}

The earlier formalization differed from Ref.~\cite{Cirac2016MPDO_arXiv} in two
linked ways. First, the canonical-form definition omitted the convention
$|\mu_1|=1$. Second, the proportional MPV comparison replaced the source's
finite power-sum identities with a convergence hypothesis on $(\mu_j^A)^N$.
Under the source normalization, this convergence requirement fails on every
sub-dominant block because the powers decay to zero.

The proof-plan memo of 2026-02-22 had already outlined a faithful finite
power-sum approach. The subsequent implementation instead introduced a
convergence-to-nonzero-limit hypothesis found in neither the source paper nor
the memo.

\paragraph{Status: supplier-side normalization.}

The supplier now implements the weight normalization stated at line~246. The
prepared weights are rescaled by their maximal modulus
\begin{align}
    m
        &= \max_k |\mu_k|,
    \label{eq:cpsv16_cf_normalization_and_proportional_comparison_supplier_maximal_modulus}
\end{align}
via weight normalization.\footnote{Formal declarations:
\leanid{MPSTensor.exists_isBNTCanonicalForm_afterBlocking_pos_normalized}
and \leanid{MPSTensor.exists_weight_normalization}.} The rescaled weights
satisfy
\begin{align}
    |\mu_k|
        &\leq 1
        \qquad\text{for every }k,
    \label{eq:cpsv16_cf_normalization_and_proportional_comparison_supplier_weight_bound}\\
    |\mu_{k_0}|
        &= 1
        \qquad\text{for some }k_0.
    \label{eq:cpsv16_cf_normalization_and_proportional_comparison_supplier_unit_witness}
\end{align}
The construction absorbs this rescaling as the per-site factor $m^N$ in the
positive-length matrix-product identity. Stated over positive lengths, this
identity matches the source formulation and involves no length-zero
coefficient.

The normalized supplier requires only that the positive-length matrix-product
family be nonzero: there exist some $N\geq 1$ and a word $w$ of length $N$
with
\begin{align}
    \tr(A^w)
        &\neq 0.
    \label{eq:cpsv16_cf_normalization_and_proportional_comparison_nonzero_positive_length_trace}
\end{align}
This condition persists under blocking. For each blocking factor $p\geq 1$,
some $q\geq 1$ satisfies
\begin{align}
    \tr((A^w)^{pq})
        &\neq 0.
    \label{eq:cpsv16_cf_normalization_and_proportional_comparison_nonzero_blocked_trace}
\end{align}
Otherwise, the Newton--Girard trace identities would imply that $(A^w)^p$,
and hence $A^w$, is nilpotent, contrary to
Eq.~\ref{eq:cpsv16_cf_normalization_and_proportional_comparison_nonzero_positive_length_trace}.
Repeating $w$ to length $pqN$ thus produces a nonzero coefficient after
blocking by $p$.

This hypothesis rules out only the identically zero positive-length family,
where the unit-modulus weight of line~246 cannot be defined. Although the
subsequent source proposition states ``for any tensor'', the convention at
line~246 presupposes a nonempty weight family. Nonzero families reflect the
source's intended scope, which the supplier hypothesis makes explicit; the
identically zero case is a degenerate reading that requires no separate
clause.

\paragraph{Status: structural realignment.}

The structural realignment is complete. The canonical-form definition now
includes the source convention $|\mu_1|=1$ as a field in the formal BNT
structure,\footnote{Formal structure: \leanid{IsNormalCanonicalFormBNT}.}
propagated by every constructor. The convergence-to-nonzero-limit fields and
the per-$N$ coefficient package have been removed. Until a source-faithful
residual comparison theorem is established, the formal boundary retains the
block-permutation and gauge-phase matching as an explicit input.

The earlier one-shot nonzero-overlap theorems and their wrappers have also been
removed. They claimed nondecay for every block from arbitrary bounded
coefficient arrays, but their hypotheses could not support that claim: a
non-dominant coefficient can vanish identically unless one encodes the source
residual peeling argument or the power-sum sector structure. The remaining
\Lean{} interfaces therefore keep the proportional matching conclusion as an
explicit hypothesis until the residual comparison theorem is proved.

The remaining task is the proof. The paper-faithful route begins with the
block-selection step of the proportional Fundamental Theorem
\cite[line~1182]{Cirac2016MPDO_arXiv}. The equal-MPV corollary then uses the
finite power-sum comparison and global gauge assembly
\cite[lines~1184--1192]{Cirac2016MPDO_arXiv}. For the dominant surviving
sector, the lower-bound argument uses
\begin{align}
    \lVert b_{N,1}\rVert
        &= 1,
    \label{eq:cpsv16_cf_normalization_and_proportional_comparison_dominant_coefficient_norm}\\
    \lVert b_{N,k}\rVert
        &\leq 1,
    \label{eq:cpsv16_cf_normalization_and_proportional_comparison_bounded_coefficient_norm}\\
    c_N
        &\neq 0.
    \label{eq:cpsv16_cf_normalization_and_proportional_comparison_nonzero_per_length_scalar}
\end{align}
The weight comparison then requires a single witness length in the strict-weight
specialization or finitely many witness lengths in the general multiplicity
setting. This development is tracked as Stage~C.\footnote{\ghissue{1559}.}

The corresponding tracked items are the canonical-form normalization, which is
structurally aligned and will close when the proofs are in place,\footnote{%
\ghissue{1554}.} and the fixed-length proportional comparison, which is
likewise structurally addressed and will close with Stage~C.\footnote{%
\ghissue{1555}.}

\bibliographystyle{alpha}
\bibliography{references}

\end{document}
